% ============================================================
% Bibliographie et nétographie
% ============================================================

\begin{thebibliography}{99}
\addcontentsline{toc}{chapter}{Bibliographie}
\markboth{\MakeUppercase{Bibliographie}}{}

\bibitem{evans}
E. Evans, \textit{Domain Driven Design: Tackling Complexity in the Heart of Software},
Addison Wesley, 2003.

\bibitem{cockburn}
A. Cockburn, \textit{Hexagonal Architecture: Ports and Adapters},
article de référence sur le patron d'architecture en ports et adaptateurs, 2005.

\bibitem{martin}
R. C. Martin, \textit{Clean Architecture: A Craftsman's Guide to Software Structure and Design},
Prentice Hall, 2017.

\bibitem{gamma}
E. Gamma, R. Helm, R. Johnson et J. Vlissides,
\textit{Design Patterns: Elements of Reusable Object Oriented Software},
Addison Wesley, 1994.

\bibitem{fowler}
M. Fowler, \textit{Patterns of Enterprise Application Architecture},
Addison Wesley, 2002.

\bibitem{vaswani}
A. Vaswani, N. Shazeer, N. Parmar, J. Uszkoreit, L. Jones, A. N. Gomez, L. Kaiser et I. Polosukhin,
\textit{Attention Is All You Need},
Advances in Neural Information Processing Systems 30, 2017.

\bibitem{lewis}
P. Lewis, E. Perez, A. Piktus, F. Petroni, V. Karpukhin, N. Goyal, H. Küttler, M. Lewis,
W. Yih, T. Rocktäschel, S. Riedel et D. Kiela,
\textit{Retrieval Augmented Generation for Knowledge Intensive NLP Tasks},
Advances in Neural Information Processing Systems 33, 2020. arXiv:2005.11401.

\bibitem{e5}
L. Wang, N. Yang, X. Huang, L. Yang, R. Majumder et F. Wei,
\textit{Multilingual E5 Text Embeddings: A Technical Report},
arXiv:2402.05672, 2024.

\bibitem{bm25}
S. E. Robertson et H. Zaragoza,
\textit{The Probabilistic Relevance Framework: BM25 and Beyond},
Foundations and Trends in Information Retrieval, vol. 3, n\textsuperscript{o} 4, p. 333 à 389, 2009.

\bibitem{rrf}
G. V. Cormack, C. L. A. Clarke et S. Buettcher,
\textit{Reciprocal Rank Fusion Outperforms Condorcet and Individual Rank Learning Methods},
Proceedings of the 32nd International ACM SIGIR Conference on Research and Development in
Information Retrieval, Boston, p. 758 à 759, 2009.

\bibitem{nogueira}
R. Nogueira et K. Cho,
\textit{Passage Re ranking with BERT},
arXiv:1901.04085, 2019.

\bibitem{mmarco}
L. Bonifacio, V. Jeronymo, H. Q. Abonizio, I. Campiotti, M. Fadaee, R. Lotufo et R. Nogueira,
\textit{mMARCO: A Multilingual Version of the MS MARCO Passage Ranking Dataset},
arXiv:2108.13897, 2021.

\bibitem{scrypt}
C. Percival et S. Josefsson,
\textit{The scrypt Password Based Key Derivation Function},
RFC 7914, Internet Engineering Task Force, 2016.

\bibitem{sha}
National Institute of Standards and Technology,
\textit{Secure Hash Standard},
FIPS PUB 180-4, 2015.

\bibitem{sip}
J. Rosenberg, H. Schulzrinne, G. Camarillo, A. Johnston, J. Peterson, R. Sparks, M. Handley
et E. Schooler,
\textit{SIP: Session Initiation Protocol},
RFC 3261, Internet Engineering Task Force, 2002.

\bibitem{webrtc}
H. Alvestrand,
\textit{Overview: Real Time Protocols for Browser Based Applications},
RFC 8825, Internet Engineering Task Force, 2021.

\bibitem{polyai}
PolyAI, \textit{Agent Studio : plateforme d'agents vocaux pour centres de contact},
documentation produit de l'éditeur, 2026.

\bibitem{cognigy}
Cognigy, \textit{Cognigy.AI : plateforme agentique pour l'expérience client},
documentation produit de l'éditeur, 2026.

\bibitem{parloa}
Parloa, \textit{AI Agent Management Platform for Contact Centers},
documentation produit de l'éditeur, 2026.

\end{thebibliography}

\newpage

% ── Nétographie ─────────────────────────────────────────────
\chapter*{Nétographie}
\addcontentsline{toc}{chapter}{Nétographie}
\markboth{\MakeUppercase{Nétographie}}{}

Les ressources suivantes constituent la documentation technique officielle des composants
employés dans le projet. Elles ont été consultées de manière continue pendant la réalisation.

\begin{description}\setlength{\itemsep}{0.4ex}

\item[{[N1]}] Documentation de LiveKit, cadriciel de communication temps réel.
\url{https://docs.livekit.io} \newline Consulté le 23 août 2026.

\item[{[N2]}] Documentation de Qdrant, moteur d'index vectoriel.
\url{https://qdrant.tech/documentation} \newline Consulté le 23 août 2026.

\item[{[N3]}] Documentation de PostgreSQL, système de gestion de base de données relationnelle.
\url{https://www.postgresql.org/docs} \newline Consulté le 23 août 2026.

\item[{[N4]}] Documentation de FastAPI, cadriciel de services web.
\url{https://fastapi.tiangolo.com} \newline Consulté le 23 août 2026.

\item[{[N5]}] Documentation de SQLAlchemy et d'Alembic, couche de persistance et gestion des
migrations. \url{https://docs.sqlalchemy.org} \newline Consulté le 23 août 2026.

\item[{[N6]}] Documentation de Deepgram, service de transcription de la parole.
\url{https://developers.deepgram.com} \newline Consulté le 23 août 2026.

\item[{[N7]}] Documentation de Cartesia, service de synthèse vocale.
\url{https://docs.cartesia.ai} \newline Consulté le 23 août 2026.

\item[{[N8]}] Documentation de l'API Gemini, service de modèle de langue.
\url{https://ai.google.dev/gemini-api/docs} \newline Consulté le 23 août 2026.

\item[{[N9]}] Spécification du Model Context Protocol, protocole de mise à disposition de
capacités. \url{https://modelcontextprotocol.io} \newline Consulté le 23 août 2026.

\item[{[N10]}] Documentation de GLPI, système de gestion de tickets, et de son interface REST.
\url{https://glpi-project.org} \newline Consulté le 23 août 2026.

\item[{[N11]}] Documentation d'OpenTelemetry, cadre de collecte de traces et de métriques.
\url{https://opentelemetry.io/docs} \newline Consulté le 23 août 2026.

\item[{[N12]}] Documentation de Prometheus, système de collecte de métriques.
\url{https://prometheus.io/docs} \newline Consulté le 23 août 2026.

\item[{[N13]}] Documentation de MinIO, stockage d'objets compatible S3.
\url{https://min.io/docs} \newline Consulté le 23 août 2026.

\item[{[N14]}] Documentation de Docker et de Docker Compose, environnement de conteneurisation.
\url{https://docs.docker.com} \newline Consulté le 23 août 2026.

\item[{[N15]}] Documentation de TanStack Start et de React, cadriciels des interfaces web.
\url{https://tanstack.com/start} \newline Consulté le 23 août 2026.

\item[{[N16]}] Fiche du modèle \texttt{intfloat/multilingual-e5-small} sur Hugging Face.
\url{https://huggingface.co/intfloat/multilingual-e5-small} \newline Consulté le 23 août 2026.

\item[{[N17]}] Fiche du modèle \texttt{cross-encoder/mmarco-mMiniLMv2-L12-H384-v1} sur
Hugging Face. \url{https://huggingface.co/cross-encoder/mmarco-mMiniLMv2-L12-H384-v1}
\newline Consulté le 23 août 2026.

\item[{[N18]}] Projet Silero VAD, détection d'activité vocale.
\url{https://github.com/snakers4/silero-vad} \newline Consulté le 23 août 2026.

\end{description}
